% Options for packages loaded elsewhere
\PassOptionsToPackage{unicode}{hyperref}
\PassOptionsToPackage{hyphens}{url}
\PassOptionsToPackage{dvipsnames,svgnames,x11names}{xcolor}
\documentclass[
  11pt,
]{article}
\usepackage{xcolor}
\usepackage{amsmath,amssymb}
\setcounter{secnumdepth}{-\maxdimen} % remove section numbering
\usepackage{iftex}
\ifPDFTeX
  \usepackage[T1]{fontenc}
  \usepackage[utf8]{inputenc}
  \usepackage{textcomp} % provide euro and other symbols
\else % if luatex or xetex
  \usepackage{unicode-math} % this also loads fontspec
  \defaultfontfeatures{Scale=MatchLowercase}
  \defaultfontfeatures[\rmfamily]{Ligatures=TeX,Scale=1}
\fi
\usepackage{lmodern}
\ifPDFTeX\else
  % xetex/luatex font selection
\fi
% Use upquote if available, for straight quotes in verbatim environments
\IfFileExists{upquote.sty}{\usepackage{upquote}}{}
\IfFileExists{microtype.sty}{% use microtype if available
  \usepackage[]{microtype}
  \UseMicrotypeSet[protrusion]{basicmath} % disable protrusion for tt fonts
}{}
\makeatletter
\@ifundefined{KOMAClassName}{% if non-KOMA class
  \IfFileExists{parskip.sty}{%
    \usepackage{parskip}
  }{% else
    \setlength{\parindent}{0pt}
    \setlength{\parskip}{6pt plus 2pt minus 1pt}}
}{% if KOMA class
  \KOMAoptions{parskip=half}}
\makeatother
\setlength{\emergencystretch}{3em} % prevent overfull lines
\providecommand{\tightlist}{%
  \setlength{\itemsep}{0pt}\setlength{\parskip}{0pt}}
% obstruction_preamble.tex — shared house style for ALL papers.
% Engine: pdfLaTeX (portable, no font cache). Injected by pandoc (see makefile).
% pandoc loads fontenc(T1)+inputenc(utf8) before this; hyperref after (so we set
% link colors via pandoc -V flags, not here).

% ---- Fonts & math ----
\IfFileExists{lmodern.sty}{\usepackage{lmodern}}{}   % Latin Modern if present
\usepackage{amsmath,amssymb,amsthm}                  % \mathbb \mathfrak \square ...
\usepackage{microtype}

% ---- Over-wide \boxed kernel statements shrink to fit the text width ----
% A long \boxed{\text{...}} line can exceed \linewidth because \text does not
% line-break in math mode. "max width=\linewidth" scales down ONLY boxes that
% are wider than the text block; short math boxes are left at natural size.
% One place, all papers.
\usepackage{adjustbox}
\let\oboldboxed\boxed
\renewcommand{\boxed}[1]{\adjustbox{max width=\linewidth}{\oboldboxed{#1}}}
\setlength{\emergencystretch}{3em}

% ---- Line breaking for verbatim and code blocks ----
\usepackage{fancyvrb}

% ---- Safe line breaking for typewriter font & long URLs/DOIs ----
\usepackage[htt]{hyphenat}
\usepackage{xurl}
\usepackage[hyphens]{url}


% ---- Flawless fix: literal Unicode (ℓ¹ Φ₁ κ¹ ⊗ →) in TEXT mode ----
% In $...$ these already work; this maps stray text-mode glyphs to math so
% nothing renders blank and you never rewrite source. If a build ever says
% "Unicode character X not set up", add one line here for that glyph.
\usepackage{newunicodechar}
\newcommand{\um}[1]{\ensuremath{#1}}
\newunicodechar{ℓ}{\um{\ell}}
\newunicodechar{Φ}{\um{\Phi}}\newunicodechar{φ}{\um{\varphi}}
\newunicodechar{Ψ}{\um{\Psi}}\newunicodechar{ψ}{\um{\psi}}
\newunicodechar{⇔}{\um{\Leftrightarrow}}
\newunicodechar{Γ}{\um{\Gamma}}\newunicodechar{γ}{\um{\gamma}}
\newunicodechar{Ω}{\um{\Omega}}\newunicodechar{ω}{\um{\omega}}
\newunicodechar{κ}{\um{\kappa}}\newunicodechar{χ}{\um{\chi}}
\newunicodechar{Λ}{\um{\Lambda}}\newunicodechar{λ}{\um{\lambda}}
\newunicodechar{Θ}{\um{\Theta}}\newunicodechar{θ}{\um{\theta}}
\newunicodechar{Δ}{\um{\Delta}}\newunicodechar{δ}{\um{\delta}}
\newunicodechar{Σ}{\um{\Sigma}}\newunicodechar{σ}{\um{\sigma}}
\newunicodechar{Π}{\um{\Pi}}\newunicodechar{π}{\um{\pi}}
\newunicodechar{α}{\um{\alpha}}\newunicodechar{β}{\um{\beta}}
\newunicodechar{ε}{\um{\varepsilon}}\newunicodechar{μ}{\um{\mu}}
\newunicodechar{ρ}{\um{\rho}}\newunicodechar{τ}{\um{\tau}}
\newunicodechar{ν}{\um{\nu}}\newunicodechar{ξ}{\um{\xi}}
\newunicodechar{ζ}{\um{\zeta}}\newunicodechar{η}{\um{\eta}}
\newunicodechar{ℝ}{\um{\mathbb{R}}}\newunicodechar{ℚ}{\um{\mathbb{Q}}}
\newunicodechar{ℤ}{\um{\mathbb{Z}}}\newunicodechar{ℕ}{\um{\mathbb{N}}}
\newunicodechar{ℂ}{\um{\mathbb{C}}}
\newunicodechar{𝒟}{\um{\mathcal{D}}}
\newunicodechar{→}{\um{\to}}\newunicodechar{↦}{\um{\mapsto}}
\newunicodechar{≤}{\um{\leq}}\newunicodechar{≥}{\um{\geq}}
\newunicodechar{≠}{\um{\neq}}\newunicodechar{≈}{\um{\approx}}
\newunicodechar{≅}{\um{\cong}}\newunicodechar{≡}{\um{\equiv}}
\newunicodechar{⊗}{\um{\otimes}}\newunicodechar{⊕}{\um{\oplus}}
\newunicodechar{∘}{\um{\circ}}\newunicodechar{∂}{\um{\partial}}
\newunicodechar{∇}{\um{\nabla}}\newunicodechar{∑}{\um{\sum}}
\newunicodechar{∏}{\um{\prod}}\newunicodechar{∞}{\um{\infty}}
\newunicodechar{⊆}{\um{\subseteq}}\newunicodechar{⊂}{\um{\subset}}
\newunicodechar{⊇}{\um{\supseteq}}\newunicodechar{∈}{\um{\in}}
\newunicodechar{∉}{\um{\notin}}\newunicodechar{∩}{\um{\cap}}
\newunicodechar{∪}{\um{\cup}}\newunicodechar{∖}{\um{\setminus}}
\newunicodechar{∅}{\um{\varnothing}}\newunicodechar{⊥}{\um{\perp}}
\newunicodechar{∀}{\um{\forall}}\newunicodechar{∃}{\um{\exists}}
\newunicodechar{±}{\um{\pm}}\newunicodechar{×}{\um{\times}}\newunicodechar{−}{\um{-}}
\newunicodechar{·}{\um{\cdot}}\newunicodechar{√}{\um{\surd}}
\newunicodechar{⟨}{\um{\langle}}\newunicodechar{⟩}{\um{\rangle}}
\newunicodechar{≪}{\um{\ll}}\newunicodechar{≫}{\um{\gg}}
\newunicodechar{∝}{\um{\propto}}\newunicodechar{≜}{\um{\triangleq}}
\newunicodechar{⇒}{\um{\Rightarrow}}\newunicodechar{⇐}{\um{\Leftarrow}}
\newunicodechar{∧}{\um{\wedge}}\newunicodechar{∨}{\um{\vee}}
\newunicodechar{⟺}{\um{\iff}}\newunicodechar{□}{\um{\square}}
\newunicodechar{⊔}{\um{\sqcup}}\newunicodechar{…}{\dots}\newunicodechar{‖}{\um{\Vert}}
\newunicodechar{▼}{\um{\blacktriangledown}}\newunicodechar{«}{\guillemotleft}\newunicodechar{»}{\guillemotright}
\newunicodechar{─}{--}\newunicodechar{│}{|}\newunicodechar{┌}{+}\newunicodechar{┐}{+}\newunicodechar{┼}{+}
\newunicodechar{═}{=}\newunicodechar{║}{\um{\parallel}}\newunicodechar{̄}{\={}}
\newunicodechar{⟹}{\um{\implies}}\newunicodechar{⟸}{\um{\impliedby}}
\newunicodechar{⁰}{\textsuperscript{0}}\newunicodechar{¹}{\textsuperscript{1}}
\newunicodechar{²}{\textsuperscript{2}}\newunicodechar{³}{\textsuperscript{3}}
\newunicodechar{⁴}{\textsuperscript{4}}\newunicodechar{⁵}{\textsuperscript{5}}
\newunicodechar{⁶}{\textsuperscript{6}}\newunicodechar{⁷}{\textsuperscript{7}}
\newunicodechar{⁸}{\textsuperscript{8}}\newunicodechar{⁹}{\textsuperscript{9}}
\newunicodechar{⁻}{\textsuperscript{-}}\newunicodechar{⁺}{\textsuperscript{+}}
\newunicodechar{ⁿ}{\textsuperscript{n}}
\newunicodechar{₀}{\textsubscript{0}}\newunicodechar{₁}{\textsubscript{1}}\newunicodechar{ᵀ}{\textsuperscript{T}}
\newunicodechar{₂}{\textsubscript{2}}\newunicodechar{₃}{\textsubscript{3}}
\newunicodechar{₄}{\textsubscript{4}}\newunicodechar{₅}{\textsubscript{5}}\newunicodechar{⋯}{\um{\cdots}}
\newunicodechar{ₜ}{\textsubscript{t}}\newunicodechar{ₑ}{\textsubscript{e}}\newunicodechar{ₙ}{\textsubscript{n}}

% ---- Geometry & spacing ----
\usepackage{geometry}\geometry{letterpaper,margin=1in}
\setlength{\parindent}{0pt}\setlength{\parskip}{0.55\baselineskip}\linespread{1.04}

% ---- Palette ----
\usepackage{xcolor}
\definecolor{oblink}{HTML}{1A4D8F}\definecolor{obrule}{HTML}{C2CCD6}
\definecolor{obboxbg}{HTML}{F6F8FA}
\definecolor{obproven}{HTML}{1B7F4B}   % PROVEN / DERIVED / STANDARD
\definecolor{obconj}{HTML}{B45309}     % CONJECTURE / OPEN / FALSE
\definecolor{obdef}{HTML}{6B21A8}      % DEFINITION / INTERPRETIVE
\definecolor{obneutral}{HTML}{475569}

% ---- Headings ----
\usepackage{titlesec}
\titleformat{\section}{\Large\bfseries}{\thesection}{0.6em}{}
\titleformat{\subsection}{\large\bfseries}{\thesubsection}{0.6em}{}
\titleformat{\subsubsection}{\normalsize\bfseries}{\thesubsubsection}{0.6em}{}
\titlespacing*{\section}{0pt}{1.4em}{0.5em}
\titlespacing*{\subsection}{0pt}{1.1em}{0.4em}

% ---- Theorem environments (for when you adopt semantic theorems) ----
\theoremstyle{plain}
\newtheorem{theorem}{Theorem}[section]
\newtheorem{lemma}[theorem]{Lemma}
\newtheorem{proposition}[theorem]{Proposition}
\newtheorem{corollary}[theorem]{Corollary}
\newtheorem{conjecture}[theorem]{Conjecture}
\theoremstyle{definition}\newtheorem{definition}[theorem]{Definition}
\theoremstyle{remark}\newtheorem*{remark}{Remark}

% ---- Blockquotes -> one uniform left-ruled callout ----
% pandoc wraps every `>` block (your theorems, Status notes, caveats) in `quote`.
\usepackage{mdframed}
\renewenvironment{quote}{%
  \begin{mdframed}[linewidth=0pt,leftline=true,linecolor=obrule,
    rightline=false,topline=false,bottomline=false,
    innerleftmargin=1em,innerrightmargin=0.8em,innertopmargin=0.5em,
    innerbottommargin=0.5em,skipabove=0.6em,skipbelow=0.6em,
    backgroundcolor=obboxbg]\small}{\end{mdframed}}

% ---- Claim-type tag macro ----
\newcommand{\claimtag}[2]{{\small\sffamily\bfseries\color{#1}[#2]}}

% ---- Inline code / file references -> standard breakable typewriter font ----
\DeclareRobustCommand{\texttt}[1]{%
  \ifmmode
    \mathtt{#1}%
  \else
    {\ttfamily\small\renewcommand{\_}{\symbol{95}\allowbreak}#1}%
  \fi}



% ---- "What this does not claim" box ----
\usepackage[most]{tcolorbox}
\newtcolorbox{nonclaim}{colback=obboxbg,colframe=obrule,boxrule=0.6pt,arc=2pt,
  left=1em,right=1em,top=0.5em,bottom=0.5em,
  title={What this does not claim},fonttitle=\bfseries,
  coltitle=black,colbacktitle=obboxbg}

% ---- Tables ----
\usepackage{booktabs}\renewcommand{\arraystretch}{1.15}
\usepackage{etoolbox}
% Compact tables, and let over-long \mathsf{CamelCase} status identifiers wrap
% inside cells. In math, \mathsf{ContinuumShadowBecomesFakeAuthority} is a single
% unbreakable atom and overruns narrow pandoc columns. Scoped to tables only, we
% re-render \mathsf as a sequence of per-word \mathsf atoms split before each
% interior capital and joined by \allowbreak, so the identifier wraps at the math
% top level (an \allowbreak between atoms IS a legal break in inline math; a text
% \mbox/\seqsplit is not). Also shrink the table font and tighten column padding.
% Body math is untouched; table cells carry only plain ASCII identifiers, so the
% splitter is emoji-safe. One place, all papers.
\let\obmathsforig\mathsf
\ExplSyntaxOn
\tl_new:N \l__obt_run_tl
\int_new:N \l__obt_i_int
\cs_new_protected:Npn \__obt_emit:
  {
    \tl_if_empty:NF \l__obt_run_tl
      {
        \int_compare:nNnT { \l__obt_i_int } > { 1 } { \allowbreak }
        \exp_args:NV \obmathsforig \l__obt_run_tl
        \int_incr:N \l__obt_i_int
        \tl_clear:N \l__obt_run_tl
      }
  }
\cs_new_protected:Npn \obtablesf #1
  {
    \tl_clear:N \l__obt_run_tl
    \int_set:Nn \l__obt_i_int { 1 }
    \tl_map_inline:nn {#1}
      {
        \str_if_eq:eeT { \str_uppercase:n {##1} } {##1}
          { \str_if_eq:eeF { \str_lowercase:n {##1} } {##1} { \__obt_emit: } }
        \tl_put_right:Nn \l__obt_run_tl {##1}
      }
    \__obt_emit:
  }
\ExplSyntaxOff
\AtBeginEnvironment{longtable}{\small\setlength{\tabcolsep}{4pt}\let\mathsf\obtablesf}
\AtBeginEnvironment{tabular}{\let\mathsf\obtablesf}

% ---- Title block polish (pandoc \maketitle) ----
% Title and subtitle both centered under each other. Body metadata stays plain
% roman (never a monospaced identity dump). Public monographs must not put
% internal programme IDs (PO-xxx, GS-xxx, …) in the title line.
\usepackage{titling}
\newcommand{\obsubtitle}{}
% Capture pandoc's \subtitle{...} (provide later is a no-op once this is defined).
\newcommand{\subtitle}[1]{\gdef\obsubtitle{#1}}
\pretitle{\begin{center}\LARGE\bfseries}
\posttitle{%
  \par
  \ifx\obsubtitle\empty\else
    \vskip0.4em
    {\large\itshape \obsubtitle\par}%
  \fi
  \end{center}\vskip0.75em}
\preauthor{\begin{center}\large}\postauthor{\par\end{center}}
\predate{\begin{center}\normalsize}\postdate{\par\end{center}\vskip0.5em\hrule\vskip1em}

% ---- Status-box markers (emoji rendered as text glyphs; one place, all papers) ----
\newunicodechar{§}{\S}
\newunicodechar{–}{--}
\newunicodechar{—}{---}
\newunicodechar{‘}{`}
\newunicodechar{’}{'}
\newunicodechar{“}{``}
\newunicodechar{”}{''}
\newunicodechar{✅}{\um{\checkmark}}
\newunicodechar{⚠}{\textbf{!}}
\newunicodechar{☣}{\um{\boxtimes}}
\newunicodechar{⊣}{\um{\dashv}}
\newunicodechar{🔷}{\textcolor{oblink}{\um{\blacklozenge}}}
\newunicodechar{️}{}
\usepackage{bookmark}
\IfFileExists{xurl.sty}{\usepackage{xurl}}{} % add URL line breaks if available
\urlstyle{same}
\hypersetup{
  pdftitle={Local Green, Global Red},
  pdfauthor={Jeremy H. Carroll},
  colorlinks=true,
  linkcolor={oblink},
  filecolor={Maroon},
  citecolor={oblink},
  urlcolor={oblink},
  pdfcreator={LaTeX via pandoc}}

\title{Local Green, Global Red}
\usepackage{etoolbox}
\makeatletter
\providecommand{\subtitle}[1]{% add subtitle to \maketitle
  \apptocmd{\@title}{\par {\large #1 \par}}{}{}
}
\makeatother
\subtitle{Detection and Recovery in Finite Linear Evaluation Models}
\author{Jeremy H. Carroll}
\date{2026-10-15}

\begin{document}
\maketitle
\begin{abstract}
Consider a finite system of sites connected by directed interfaces. Each
interface carries a non-expansive linear translation, and each evaluator
observes a subset of interface defects through non-expansive linear
measurements. We establish four principal conclusions. First, every
evaluator score is bounded above by the total defect mass. Second, all
evaluators score zero exactly when every interface defect lies in the
joint kernel of that interface's watchers. Third, complete interface
coverage does not ensure soundness: an explicit construction with three
component sites, three overlap sites, six interfaces, and three
evaluators has an identically zero score vector while its total defect
mass can equal any prescribed \(R > 0\). Consequently, no finite-valued
function of the maximum evaluator score can uniformly upper-bound the
true defect on this system. Fourth, trivial joint kernels at every
interface are sufficient but not necessary for soundness, because
coupling among realizable interface defects can eliminate otherwise
invisible directions. The finite-dimensional quantitative recovery bound
is proved in prose by compactness. Lean 4 formalizes the specified
abstract observer statements and edge-level identities; rational repair
identities include squared Euclidean quantities, while the real infima
and square roots are established by the written proofs. The results
apply only when a concrete evaluation system instantiates the stated
finite linear model.
\end{abstract}

\textbf{Series:}
\texttt{001\_finite\_obstruction\_calculus\ /\ 001\_local\_green\_global\_red\_foundational\_series\ /\ lggr\_v\_local\_green\_global\_red}\\
\textbf{Paper ID:} \texttt{lggr\_v\_local\_green\_global\_red}\\
\textbf{Date:} 2026-10-15\\
\textbf{Version:} 1.3.2\\
\textbf{Claim Ledger:}
\texttt{lggr\_v\_local\_green\_global\_red\_claim\_ledger.md}\\
\textbf{Change Log:}
\texttt{lggr\_v\_local\_green\_global\_red\_change\_log.md}\\
\textbf{License:} CC-BY-4.0

\begin{center}\rule{0.5\linewidth}{0.5pt}\end{center}

\subsection{Why this matters}\label{why-this-matters}

A dashboard can test every interface and still miss inconsistency if its
measurements discard the relevant directions. This note makes that
observation precise in a finite linear model. A fixed covered system
admits assignments with the same all-zero score vector and arbitrarily
large hidden defect mass.

Soundness depends on whether any nonzero realizable defect is invisible
to all measurements. The note also gives a coupled system in which
individually incomplete measurements jointly rule out hidden defects.
Applications must specify the linear model and measurements to which
these results are being applied.

\subsection{1. The objects}\label{the-objects}

Let \(X\) and \(E\) be finite sets of sites and interfaces, and let
\(A\) be a finite nonempty set of evaluators. Here ``finite'' refers to
these index sets. The state spaces are real normed vector spaces; they
need not be finite-dimensional unless explicitly stated. Assignments are
unrestricted unless a state bound is specified.

\begin{itemize}
\tightlist
\item
  A finite set of \textbf{sites} and a finite set of directed
  \textbf{interfaces} between them. Each site \(x\) carries a real
  normed vector space \(F(x)\) (its state space), and each interface
  \(e\) from site \(a\) to site \(b\) carries a linear map
  \(r_e : F(b)\to F(a)\) with \(\lVert r_e v\rVert\le\lVert v\rVert\) (a
  lossy translation of \(b\)'s state into \(a\)'s terms).
\item
  An \textbf{assignment} \(s\) gives each site a state \(s(x)\in F(x)\).
  The \textbf{defect} of \(s\) at interface \(e=(a,b)\) is
  \[\Delta(s)_e := r_e\,s(b) - s(a) \in F(a),\] the disagreement between
  what \(b\) claims (translated) and what \(a\) holds. The
  \textbf{global inconsistency} is the total defect mass
  \[I(s) := \sum_e \lVert\Delta(s)_e\rVert .\] \(I(s)=0\) means the
  assignment is globally consistent (``Global Green''); \(I(s)>0\) is
  ``Global Red''. If site states are bounded by
  \(\lVert s(x)\rVert \le B\), then because translations are contractive
  (\(\lVert r_e v\rVert\le\lVert v\rVert\)), each defect satisfies
  \(\lVert\Delta(s)_e\rVert \le \lVert r_e s(b)\rVert + \lVert s(a)\rVert \le 2B\),
  yielding the contractivity bound \[I(s) \le 2B |E|.\] In particular,
  producing arbitrarily large global inconsistency requires an unbounded
  family of assignments.
\item
  An evaluator \(\alpha \in A\) is a pair: a set
  \(E^\alpha \subseteq E\) of interfaces it watches, and for each
  watched interface a linear measurement
  \(q_e^\alpha : F(a)\to V_e^\alpha\) into some normed space, with
  \(\lVert q_e^\alpha v\rVert\le\lVert v\rVert\) (measurements do not
  amplify). Its \textbf{score} is
  \[J^\alpha(s) := \sum_{e\in E^\alpha}\bigl\lVert q_e^\alpha(\Delta(s)_e)\bigr\rVert .\]
  \(J^\alpha(s)=0\) is ``Local Green'' for \(\alpha\); the family is
  \textbf{All-Green} when every evaluator scores zero.
\item
  The \textbf{blind spot} of interface \(e\) is the set \(K_e\) of
  defect values that \emph{every} watcher of \(e\) measures as zero (the
  intersection of the measurements' kernels). An unwatched interface's
  blind spot is everything.
\end{itemize}

That is the entire setup. An application of this note must name its
sites, interfaces, state spaces, translation maps, evaluators, and
measurements; if it cannot, the theorems below say nothing about it.

\subsection{2. Monitoring is one-way}\label{monitoring-is-one-way}

\textbf{Theorem 1 (evaluators cannot create inconsistency).} For every
assignment \(s\) and every evaluator \(\alpha\):
\(J^\alpha(s)\le I(s)\). In particular, a globally consistent system
makes every evaluator report green.

\emph{Proof sketch.} Each measured term is at most the true term because
measurements do not amplify; then a sum over some interfaces is at most
the sum over all. \(\blacksquare\)

The useful direction would be the converse --- do green evaluators
certify a consistent system? Everything below is about exactly when that
converse holds.

\subsection{3. The exact mechanism of false
green}\label{the-exact-mechanism-of-false-green}

\textbf{Theorem 2 (All-Green, characterized).} The family is All-Green
at \(s\) \textbf{if and only if} at every interface, the defect lies in
that interface's blind spot: \(\Delta(s)_e\in K_e\) for all \(e\).

\emph{Proof.} Each score is a finite sum of nonnegative terms, so it
vanishes exactly when every measured term vanishes. Since measurement
codomains are normed spaces, a measured term has zero norm exactly when
its measured vector is zero. Requiring this for every evaluator is
precisely \(\Delta(s)_e\in K_e\) for every interface. The unwatched case
imposes no condition. \(\blacksquare\)

So ``all evaluators green'' does not, without further assumptions, imply
no defect; it means ``no defect \emph{outside the blind spots}''.
Whether that is reassuring depends entirely on the blind spots, which
brings us to:

\textbf{Theorem 3 (false green with full coverage, at any severity).}
There is an explicit system --- three component sites \(u_0, u_1, u_2\)
and three overlap sites \(x_{01}, x_{12}, x_{20}\) (six sites total)
connected by six directed interfaces \(e_{01}^0 = (x_{01}, u_0)\),
\(e_{01}^1 = (x_{01}, u_1)\), \(e_{12}^1 = (x_{12}, u_1)\),
\(e_{12}^2 = (x_{12}, u_2)\), \(e_{20}^2 = (x_{20}, u_2)\), and
\(e_{20}^0 = (x_{20}, u_0)\), state spaces \(\mathbb R^3\) with the
\(\ell^1\) norm, identity translations \(r_e = I_3\) --- and three
evaluators \(\alpha \in \{0, 1, 2\}\) that jointly watch \emph{every}
interface, such that for every number \(R>0\) there is an assignment
\(s_R\) with:

\begin{itemize}
\tightlist
\item
  every evaluator scores exactly zero (All-Green), and
\item
  \(I(s_R) = R\).
\end{itemize}

The six-site construction is explicit: for any amplitude
\(\lambda \in \mathbb R\), define the assignment \(s^{(\lambda)}\) by
setting \(s^{(\lambda)}(u_i) = \lambda e_i\) for each component site
\(i \in \{0, 1, 2\}\), where \(e_i\) is the standard basis vector in
\(\mathbb R^3\), and \(s^{(\lambda)}(x_{ij}) = 0\) on all overlap sites.
On each interface into component site \(u_i\), the defect is
\(\Delta(s^{(\lambda)})_e = r_e s^{(\lambda)}(u_i) - s^{(\lambda)}(x_{ij}) = \lambda e_i\).
Evaluator \(i \in \{0, 1, 2\}\) watches the two interfaces incident to
component site \(u_i\) via the measurement \(q^i\) that deletes
coordinate \(i\). Because the defects at those interfaces point
precisely along coordinate \(i\), each evaluator measures zero:
\(J^i(s^{(\lambda)}) = 0\). The global inconsistency is
\[I(s^{(\lambda)}) = \sum_{e} \lVert\Delta(s^{(\lambda)})_e\rVert_1 = 6 \lVert\lambda e_i\rVert_1 = 6 |\lambda|.\]
Setting \(s_R := s^{(R/6)}\) gives the severity-indexed assignment with
\(I(s_R) = 6(R/6) = R\) and All-Green score vector \((0, 0, 0)\).

\emph{Why this construction uses overlap sites (not a minimality
claim).} On a literal bidirected three-site triangle without overlap
sites, interfaces between sites \(u\) and \(v\) in opposite directions
have defects \(s(v) - s(u) = -(s(u) - s(v))\). Opposite defects are
negatives of each other; requiring them to lie along different
coordinate axes forces them to vanish. Introducing the three overlap
sites \(x_{ij}\) with \(s(x_{ij}) = 0\) decouples the interfaces.
Setting the overlap states to zero makes both interfaces at component
\(u_i\) carry the same defect \(\lambda e_i\). This realizes the
required hidden directions without the opposite-edge constraint of the
literal bidirected triangle; it does not make all six defects
independent. Six sites are used for this construction, with no
minimality claim.

\textbf{Theorem 4 (no recovery function).} Over the fixed system of
Theorem 3, there is \textbf{no finite-valued function}
\(f : \mathbb{R}_{\ge 0} \to \mathbb{R}_{\ge 0}\) with
\[I(s)\ \le\ f\Bigl(\max_\alpha J^\alpha(s)\Bigr)\quad\text{for all assignments }s ;\]
and in the stronger \textbf{joint} form, no finite-valued function
\(F : \mathbb{R}_{\ge 0}^3 \to \mathbb{R}_{\ge 0}\) of the full score
vector with
\[I(s)\ \le\ F\bigl(J^0(s),\,J^1(s),\,J^2(s)\bigr)\quad\text{for all assignments }s .\]
The finite-valued qualification excludes the vacuous extended-real bound
\(F\equiv+\infty\). Every fixed finite rescaling of the zero score
vector still leaves that vector compatible with arbitrarily large hidden
defect mass. This conclusion uses unrestricted assignments; on a
state-bounded family, the bound \(I(s)\le 2B|E|\) from §1 remains
available. Aggregating, weighting, or post-processing that vector cannot
manufacture a finite uniform upper bound.

\emph{Proof.} Suppose such a function \(f\) (or \(F\)) exists.
Evaluating the bound on the family of assignments \(s_R\) from Theorem 3
with \(R > 0\), we have \(J^\alpha(s_R) = 0\) for all
\(\alpha \in \{0, 1, 2\}\), so \(\max_\alpha J^\alpha(s_R) = 0\) and the
score vector is \((0, 0, 0)\). The bound requires
\(R = I(s_R) \le f(0)\) (or \(R \le F(0, 0, 0)\)) for every \(R > 0\).
Taking \(R = f(0) + 1\) (or \(F(0, 0, 0) + 1\)) produces an immediate
contradiction since \(f(0) < \infty\). \(\blacksquare\)

\subsection{4. What actually restores soundness --- and what does
not}\label{what-actually-restores-soundness-and-what-does-not}

For an interface \(e=(a(e),b(e))\), define the assignment and defect
spaces by \[C^0:=\prod_{x\in X}F(x),\qquad
C^1:=\prod_{e\in E}F(a(e)).\] The linear defect operator
\(D: C^0\to C^1\) is \[(Ds)_e=r_e s(b(e))-s(a(e)).\] Equip \(C^1\) with
the sum norm
\[\lVert\omega\rVert_{1,E}:=\sum_{e\in E}\lVert\omega_e\rVert_{F(a(e))}.\]
Thus \(I(s)=\lVert Ds\rVert_{1,E}\). Write \(K:=\prod_{e\in E}K_e\).

\textbf{Soundness Criterion.} All-Green implies \(I(s)=0\) for every
assignment \(s\in C^0\) if and only if
\[\operatorname{im}D\cap K=\{0\}.\]

\emph{Proof.} By Theorem 2, All-Green is equivalent to \(Ds\in K\).
Every such defect also belongs to \(\operatorname{im}D\), and \(I(s)=0\)
exactly when \(Ds=0\). Thus soundness says that
\(\operatorname{im}D\cap K\) contains no nonzero vector. Both spaces
contain zero, giving the stated equality. \(\blacksquare\)

The criterion concerns realizable defects. The sufficient condition in
Theorem 5 makes all ambient blind spots trivial; Theorem 6 shows how
coupling can give soundness without that stronger condition. The
relation to the formalization's seminormed hypotheses is recorded in §6.

Let \(Q\) stack all the individual linear measurements of defect
cochains. Then \(\ker Q=K\), and the soundness criterion is injectivity
of \(Q\) on \(\operatorname{im}D\). This is the familiar null-space
criterion for linear observation; compare {[}2{]}. The present note
expresses it for composed interface evaluation, gives explicit covered
and coupled examples, and supplies accompanying formal proof artifacts.

\textbf{Theorem 5 (a sufficient condition).} If every interface is
watched and every interface's blind spot is trivial (\(K_e=\{0\}\) ---
the watchers of each interface \emph{jointly} separate all states,
though no single watcher need do so), then All-Green implies \(I(s)=0\).

This corresponds to the special case where the ambient product
\(\prod_e K_e = \{0\}\) is trivial, so its intersection with
\(\operatorname{im} D\) is automatically \(\{0\}\) regardless of the
system's coupling.

\textbf{Theorem 6 (the condition is not necessary).} Trivial blind spots
are not required for soundness. There is a second explicit system ---
one shared site under two others, with \emph{zero maps} as translations,
so that both interface defects are forced to be the same vector
\(v \in \mathbb R^2\) --- monitored by two evaluators with coordinate
projections \(q^0(v) = v_1\) and \(q^1(v) = v_2\) (whose blind spots are
the coordinate axes \(K_0 = \{0\} \times \mathbb R\) and
\(K_1 = \mathbb R \times \{0\}\)). Each evaluator has a full line as its
blind spot, yet All-Green still forces \(I(s)=0\): the shared defect
\(v\) must satisfy \(q^0(v) = v_1 = 0\) and \(q^1(v) = v_2 = 0\),
forcing \(v = (0, 0) = 0\). Here \(\prod_e K_e \ne \{0\}\), but
structural \emph{coupling} ensures
\(\operatorname{im} D \cap \prod_e K_e = \{0\}\). Coupling between
interfaces can compensate for individual blindness.

\emph{Proof.} Let every site carry \(\mathbb R^2\) with its \(\ell^1\)
norm. Take sites \(\{a_0,a_1,b\}\), interfaces \(e_i=(b,a_i)\), and zero
translations. For every assignment \(s\), both interface defects equal
\(-s(b)\), independently of \(s(a_0)\) and \(s(a_1)\). Evaluator 0
watches \(e_0\) and measures the first coordinate; evaluator 1 watches
\(e_1\) and measures the second. Their blind spots are respectively
\(\{0\}\times\mathbb R\) and \(\mathbb R\times\{0\}\), so both are
nontrivial and every interface is watched. If both scores vanish, both
coordinates of \(s(b)\) vanish. Hence both defects vanish and
\(I(s)=0\). This demonstrates soundness despite nontrivial blind spots
on both interfaces. \(\blacksquare\)

The design consequence cuts both ways. A family failing the Theorem 5
condition is not automatically unsound --- but proving it sound then
requires an argument about the specific coupling, and Theorem 3 shows
what happens when nobody makes that argument.

\textbf{Remark (Raw inconsistency and distance to the repair image).}
The raw inconsistency is \(I(s)=\lVert Ds\rVert_{1,E}\). For an
arbitrary defect cochain \(\omega\in C^1\), define its distance to the
image of the allowed repair operator by
\[\Phi(\omega):=\inf_{t\in C^0}\lVert\omega-Dt\rVert_{1,E},\] the
infimal residual under adjustments of site states \(t\). Exactness means
\(\omega\in\operatorname{im}D\) and implies \(\Phi(\omega)=0\). The
converse requires closedness of the image in the chosen norm; it holds
for the finite-dimensional diagonal image in the example below but is
not asserted for arbitrary normed fibers. A positive value of \(\Phi\)
excludes even approximation by repairs.

\textbf{Proposition (finite-dimensional quantitative recovery; proved in
prose).} If the realizable defect space \(Y=\operatorname{im}D\) is
finite-dimensional and soundness holds, there is a constant \(C>0\),
depending on the fixed system and measurements, such that
\[I(s)\le C\max_{\alpha\in A}J^\alpha(s)\quad\text{for all }s\in C^0.\]
\emph{Proof.} Put
\[p(\omega):=\max_{\alpha\in A}\sum_{e\in E^\alpha}\lVert q_e^\alpha(\omega_e)\rVert.\]
If \(Y\ne\{0\}\), the continuous function \(p\) is positive on the
compact unit sphere of \(Y\), so its minimum \(m\) there is positive.
Homogeneity gives \(\lVert\omega\rVert_{1,E}\le p(\omega)/m\) on \(Y\);
take \(C=1/m\). If \(Y=\{0\}\), the bound is immediate, for example with
\(C=1\). \(\blacksquare\)

This is a proved finite-dimensional result, not Lean-mechanized in this
deposit. The expression \(\inf_{s:I(s)=1}\max_\alpha J^\alpha(s)\) is
used only when \(Y\ne\{0\}\); otherwise its constraint set is empty. The
floating-point helper \texttt{compute\_recovery\_constant} instead uses
ambient edgewise injectivity, singular values, and norm-conversion
factors. It need not detect recovery supplied only by coupling and does
not compute this compact-sphere optimum.

Crucially, these represent two entirely distinct failure modes:

\begin{enumerate}
\def\labelenumi{\arabic{enumi}.}
\tightlist
\item
  \textbf{Raw Inconsistency / False Green:} The three-patch witness of
  Theorem 3 has \(I(s_R) = R > 0\), but because its defect is generated
  by an actual section \(s_R\), \(\Delta(s_R) \in \operatorname{im} D\),
  and hence \(\Phi(\Delta(s_R)) = 0\) by definition.
\item
  \textbf{Quotient Inconsistency / Non-Exact Defect:} On the coupled
  system of Theorem 6 with scalar states, the defect operator
  \(Ds = (-\alpha, -\alpha)\) has diagonal image
  \(\operatorname{im} D = \{(-\alpha, -\alpha) : \alpha \in \mathbb R\}\).
  The test cochain \(\omega = (1, -1)\) cannot arise as \(\Delta(s)\)
  for any section at all: \(\omega \notin \operatorname{im} D\). For
  this cochain,
  \[\Phi(\omega)\ =\ \inf_{\alpha \in \mathbb R} \bigl(|1 + \alpha| + |-1 + \alpha|\bigr)\ =\ 2\ >\ 0,\]
  Indeed, for every real \(\alpha\),
  \[|1+\alpha|+|\alpha-1|\ge |(1+\alpha)-(\alpha-1)|=2,\] and
  \(\alpha=0\) attains two (machine-checked by
  \texttt{repair\_mass\_ge\_two}, \texttt{repair\_mass\_attained}, and
  \texttt{omega\_not\_exact}).
\end{enumerate}

\textbf{Lemma (Infimal Minkowski repair distances).} Let the scalar
coupled repair image be \(\{(-\alpha,-\alpha):\alpha\in\mathbb R\}\) and
let \(\omega=(1,-1)\). The residual is
\(r(\alpha)=(1+\alpha,\alpha-1)\). For \(p\in\{1,2,\infty\}\) set
\[\Phi_p(\omega)=\inf_{\alpha\in\mathbb R}\|r(\alpha)\|_p.\] Then
\(\Phi_1(\omega)=2\), \(\Phi_2(\omega)=\sqrt{2}\), and
\(\Phi_\infty(\omega)=1\). The value \(2\) is attained at every
\(\alpha\in[-1,1]\). The values \(\sqrt{2}\) and \(1\) are attained only
at \(\alpha=0\).

\emph{Proof.} For \(\alpha\in[-1,1]\) one has \(1+\alpha\ge 0\) and
\(1-\alpha\ge 0\), so \[\|r(\alpha)\|_1=(1+\alpha)+(1-\alpha)=2.\] For
\(\alpha>1\) the sum is \(2\alpha>2\). For \(\alpha<-1\) the sum is
\(-2\alpha>2\). Thus the infimum is \(2\), attained exactly on
\([-1,1]\). The inequality \(\ge 2\) is the comparison already used for
\(\Phi\); the attainment set is new.

Next,
\[\|r(\alpha)\|_2^2=(1+\alpha)^2+(\alpha-1)^2=2\alpha^2+2=2(\alpha^2+1)\ge 2,\]
with equality if and only if \(\alpha=0\). The square root is increasing
on \([0,\infty)\), so
\[\|r(\alpha)\|_2=\sqrt{2}\,\sqrt{\alpha^2+1}\ge\sqrt{2},\] with
equality only at \(\alpha=0\).

If \(\alpha\ge 0\), then
\(\|r(\alpha)\|_\infty\ge|1+\alpha|=1+\alpha\ge 1\), so equality forces
\(\alpha=0\). If \(\alpha\le 0\), then
\(\|r(\alpha)\|_\infty\ge|\alpha-1|=1-\alpha\ge 1\), so equality forces
\(\alpha=0\). At \(\alpha=0\), \(\|r(0)\|_\infty=\max(1,1)=1\).

The chain \(1\le\sqrt{2}\le 2\) is therefore the chain of these three
infima. Evaluating the three norms only at \(\alpha=0\) gives upper
bounds on the infima. The matching lower bounds are the three
calculations above. The \(\ell^1\) identity and the squared \(\ell^2\)
identity are rational. The symbol \(\sqrt{2}\) names the positive square
root of the proved square \(2\); the real square-root conclusion is
proved above and has no separate Lean declaration here. \(\blacksquare\)

The lemma is one cochain and one diagonal repair line. The next lemma
keeps that image and replaces the single residual by the residual of an
arbitrary scalar pair. Neither lemma computes a quotient distance for a
general norm or a non-diagonal image. The rational content --- the
\(\ell^1\) attainment interval, the squared Euclidean identity and its
unique minimizer, and the \(\ell^\infty\) identity --- is
\texttt{MinkowskiRepair.r2\_l1\_iff},
\texttt{MinkowskiRepair.r2\_l2\_sq\_iff}, and
\texttt{MinkowskiRepair.r2\_linf\_iff} in
\texttt{lean4/minkowski\_repair.lean}. That file is not a Lake root.
\texttt{quotient\_attained\_minkowski\_chain} still evaluates the three
norms only at \(\alpha=0\) and does not prove these infima. The positive
square roots name the proved squares. The pointwise comparison
\(\|r(\alpha)\|_\infty\le\|r(\alpha)\|_2\le\|r(\alpha)\|_1\) does not by
itself name the infima. No dashboard margin, sheaf, or observation map
is added.

\textbf{Lemma (Minkowski distances of a stated scalar cochain).} Let the
scalar coupled repair image remain
\(\{(-\alpha,-\alpha):\alpha\in\mathbb R\}\). Let \(\omega=(a,b)\) with
\(a,b\in\mathbb R\), and let the residual be
\(r(\alpha)=(a+\alpha,\,b+\alpha)\). For \(p\in\{1,2,\infty\}\) set
\[\Phi_p(\omega)=\inf_{\alpha\in\mathbb R}\|r(\alpha)\|_p.\] Then
\[\Phi_1(\omega)=|a-b|,\qquad\Phi_2(\omega)=\frac{|a-b|}{\sqrt{2}},\qquad\Phi_\infty(\omega)=\frac{|a-b|}{2}.\]
Let \(m=\min(a,b)\) and \(M=\max(a,b)\). The value \(|a-b|\) is attained
at every \(\alpha\in[-M,-m]\). The values \(|a-b|/\sqrt{2}\) and
\(|a-b|/2\) are attained only at \(\alpha=-(a+b)/2\). In particular
\(\Phi_p(\omega)=0\) if and only if \(a=b\), in which case the only
attaining repair is \(\alpha=-a\) and \(\omega\) lies on the repair
image. The pair \((a,b)=(1,-1)\) is the preceding lemma.

Evaluating the three norms at the zero repair \(\alpha=0\) produces
\(\|(a,b)\|_p\). Each such number is only an upper bound on
\(\Phi_p(\omega)\). It equals \(\Phi_1(\omega)\) if and only if
\(ab\le 0\), and it equals \(\Phi_2(\omega)\) or \(\Phi_\infty(\omega)\)
if and only if \(b=-a\).

\emph{Proof.} For the \(\ell^1\) norm and every real \(\alpha\),
\[|a+\alpha|+|b+\alpha|\ge|(a+\alpha)-(b+\alpha)|=|a-b|.\] Set
\(t=-\alpha\). If \(m\le t\le M\), then \(t\) lies between \(a\) and
\(b\), so the sum equals \(|a-b|\). If \(t>M\), then \(t>a\) and
\(t>b\), and \[|a-t|+|b-t|=2t-(a+b)=|a-b|+2(t-M),\] using \(a+b=m+M\)
and \(|a-b|=M-m\). If \(t<m\), the same expansion gives
\(|a-b|+2(m-t)\). Thus the infimum is \(|a-b|\), attained exactly for
\(t\in[m,M]\), that is for \(\alpha\in[-M,-m]\).

For the Euclidean norm,
\[\|r(\alpha)\|_2^2=(a+\alpha)^2+(b+\alpha)^2=2\alpha^2+2(a+b)\alpha+a^2+b^2=2\left(\alpha+\frac{a+b}{2}\right)^2-\frac{(a+b)^2}{2}+a^2+b^2.\]
The constant term equals \((a-b)^2/2\), because
\[-\frac{(a+b)^2}{2}+a^2+b^2=\frac{-(a^2+2ab+b^2)+2a^2+2b^2}{2}=\frac{(a-b)^2}{2}.\]
Hence
\[\|r(\alpha)\|_2^2=2\left(\alpha+\frac{a+b}{2}\right)^2+\frac{(a-b)^2}{2}.\]
The squared shift is nonnegative and vanishes only at
\(\alpha=-(a+b)/2\). The square root is increasing on \([0,\infty)\), so
\[\|r(\alpha)\|_2\ge\frac{|a-b|}{\sqrt{2}},\] with equality only at that
repair.

For the supremum norm, set \(\delta=|a-b|/2\). Then
\[\max(|a-t|,|b-t|)\ge\frac{|a-t|+|b-t|}{2}\ge\delta.\] At \(t=(a+b)/2\)
both absolute values equal \(\delta\), so the maximum equals \(\delta\).
Suppose the maximum equals \(\delta\). Then \(t\) lies in both
\([a-\delta,a+\delta]\) and \([b-\delta,b+\delta]\). If \(a\le b\),
these intervals are \([a-\delta,(a+b)/2]\) and \([(a+b)/2,b+\delta]\),
whose intersection is the singleton \(\{(a+b)/2\}\). The case \(b\le a\)
is symmetric. Hence \(\Phi_\infty(\omega)=\delta=|a-b|/2\), attained
only at \(\alpha=-(a+b)/2\).

The three formulas give the chain
\(\Phi_\infty(\omega)\le\Phi_2(\omega)\le\Phi_1(\omega)\). They vanish
simultaneously exactly when \(a=b\).

The zero-repair residual is \(r(0)=(a,b)\), so
\(\|(a,b)\|_p\ge\Phi_p(\omega)\). For \(p=1\), the repair \(\alpha=0\)
lies in \([-M,-m]\) if and only if \(m\le 0\le M\), equivalently
\(ab\le 0\); then \(|a|+|b|=|a-b|\). If \(ab>0\), then \(\alpha=0\) lies
outside that interval and \(|a|+|b|>|a-b|\). For \(p\in\{2,\infty\}\)
the only attaining repair is \(\alpha=0\) if and only if \(a+b=0\). If
\(a+b\ne 0\), the zero-repair value is strictly larger than the infimum.
\(\blacksquare\)

The lemma is every cochain in \(\mathbb R^2\) and the same diagonal
repair line. It does not compute a quotient distance for a norm other
than \(p\in\{1,2,\infty\}\), a non-diagonal image, a third component, or
a vector-valued fiber. The rational content is
\texttt{MinkowskiRepair.diag\_l1\_iff},
\texttt{MinkowskiRepair.diag\_l2\_sq\_iff}, and
\texttt{MinkowskiRepair.diag\_linf\_iff} in
\texttt{lean4/minkowski\_repair.lean}. The symbol \(|a-b|/\sqrt{2}\)
names the positive square root of the proved square \((a-b)^2/2\). The
declaration \texttt{quotient\_attained\_minkowski\_chain} still
evaluates a residual only at \(\alpha=0\), and for a general pair that
evaluation is the upper bound above, not the infimum. No dashboard
margin, sheaf, or observation map is added. The next lemma leaves this
diagonal image and computes the same three distances on one stated
non-diagonal line.

\textbf{Lemma (Minkowski distances on one non-diagonal line).} Let the
repair image be \(\{(-\alpha,-2\alpha):\alpha\in\mathbb R\}\), the line
through the origin in the direction \((-1,-2)\). It is not the diagonal.
Let \(\omega=(a,b)\) with \(a,b\in\mathbb R\), and let the residual be
\(r(\alpha)=(a+\alpha,\,b+2\alpha)\). For \(p\in\{1,2,\infty\}\) set
\[\Phi_p(\omega)=\inf_{\alpha\in\mathbb R}\|r(\alpha)\|_p.\] Then
\[\Phi_1(\omega)=\frac{|2a-b|}{2},\qquad\Phi_2(\omega)=\frac{|2a-b|}{\sqrt{5}},\qquad\Phi_\infty(\omega)=\frac{|2a-b|}{3}.\]
The attainment sets are the singletons
\[A_1(\omega)=\left\{-\frac{b}{2}\right\},\qquad A_2(\omega)=\left\{-\frac{a+2b}{5}\right\},\qquad A_\infty(\omega)=\left\{-\frac{a+b}{3}\right\}.\]
Thus \(\Phi_1\) is attained only at \(\alpha=-b/2\), \(\Phi_2\) only at
\(\alpha=-(a+2b)/5\), and \(\Phi_\infty\) only at \(\alpha=-(a+b)/3\).
The three distances vanish if and only if \(b=2a\), in which case
\(\omega\) lies on the image and the only attaining repair is
\(\alpha=-a\). The three attainment sets coincide if and only if
\(b=2a\). Since \(2<\sqrt{5}<3\), the coefficients satisfy
\(1/3<1/\sqrt{5}<1/2\), so
\(\Phi_\infty(\omega)\le\Phi_2(\omega)\le\Phi_1(\omega)\), and any two
of the three distances agree if and only if \(b=2a\). The pair
\((a,b)=(1,-1)\) gives \(\Phi_1=3/2\), \(\Phi_2=3/\sqrt{5}\), and
\(\Phi_\infty=1\), with attainment sets \(\{1/2\}\), \(\{1/5\}\), and
\(\{0\}\). Those are not the diagonal values of the preceding lemmas.

Evaluating the three norms at the zero repair \(\alpha=0\) produces
\(\|(a,b)\|_p\). Each such number is only an upper bound on
\(\Phi_p(\omega)\). It equals \(\Phi_1(\omega)\) if and only if \(b=0\),
it equals \(\Phi_2(\omega)\) if and only if \(a+2b=0\), and it equals
\(\Phi_\infty(\omega)\) if and only if \(a+b=0\).

\emph{Proof.} For every real \(\alpha\),
\[|2a-b|=|2(a+\alpha)-(b+2\alpha)|\le 2|a+\alpha|+|b+2\alpha|\le 2\bigl(|a+\alpha|+|b+2\alpha|\bigr).\]
The second inequality is an equality if and only if \(b+2\alpha=0\).
When \(b+2\alpha=0\), the first inequality is an equality for every
\(a\). Thus \[\|r(\alpha)\|_1\ge\frac{|2a-b|}{2},\] with equality only
at \(\alpha=-b/2\). At that repair the residual is \((a-b/2,\,0)\),
whose \(\ell^1\) norm is \(|2a-b|/2\).

For the Euclidean norm,
\[\|r(\alpha)\|_2^2=(a+\alpha)^2+(b+2\alpha)^2=5\alpha^2+2(a+2b)\alpha+a^2+b^2=5\left(\alpha+\frac{a+2b}{5}\right)^2-\frac{(a+2b)^2}{5}+a^2+b^2.\]
The constant term equals \((2a-b)^2/5\), because
\[-\frac{(a+2b)^2}{5}+a^2+b^2=\frac{-(a^2+4ab+4b^2)+5a^2+5b^2}{5}=\frac{4a^2-4ab+b^2}{5}=\frac{(2a-b)^2}{5}.\]
Hence
\[\|r(\alpha)\|_2^2=5\left(\alpha+\frac{a+2b}{5}\right)^2+\frac{(2a-b)^2}{5}.\]
The squared shift is nonnegative and vanishes only at
\(\alpha=-(a+2b)/5\). The square root is increasing on \([0,\infty)\),
so \[\|r(\alpha)\|_2\ge\frac{|2a-b|}{\sqrt{5}},\] with equality only at
that repair.

For the supremum norm, set \(\delta=|2a-b|/3\). The same linear
combination gives
\[|2a-b|\le 2|a+\alpha|+|b+2\alpha|\le 3\max\bigl(|a+\alpha|,|b+2\alpha|\bigr),\]
so \(\|r(\alpha)\|_\infty\ge\delta\). At \(\alpha=-(a+b)/3\),
\[a+\alpha=\frac{2a-b}{3},\qquad b+2\alpha=-\frac{2a-b}{3},\] and the
maximum equals \(\delta\). Suppose the maximum equals \(\delta\). Then
both inequalities above are equalities, so
\(|a+\alpha|=|b+2\alpha|=\delta\) and \(|2x+y|=2|x|+|y|\) for
\(x=a+\alpha\) and \(y=-(b+2\alpha)\). For real numbers,
\(|u+v|=|u|+|v|\) if and only if \(uv\ge 0\). Taking \(u=2x\) and
\(v=y\) gives \(xy\ge 0\), hence \((a+\alpha)(b+2\alpha)\le 0\). Two
real numbers of equal absolute value whose product is at most zero are
negatives of each other, so \(a+\alpha=-(b+2\alpha)\), that is
\(\alpha=-(a+b)/3\). Thus \(\Phi_\infty(\omega)=\delta\), attained only
at that repair.

Since \(2<\sqrt{5}<3\), the coefficients \(1/3\), \(1/\sqrt{5}\), and
\(1/2\) are strictly ordered. The three formulas therefore give
\(\Phi_\infty(\omega)\le\Phi_2(\omega)\le\Phi_1(\omega)\), and any two
of the distances agree if and only if \(b=2a\). They vanish
simultaneously exactly when \(b=2a\). In that case each stated
attainment point equals \(-a\), and \(r(-a)=(0,0)\). Conversely, if the
three singletons are equal, then \(-b/2=-(a+b)/3\), so \(b=2a\).

The zero-repair residual is \(r(0)=(a,b)\), so
\(\|(a,b)\|_p\ge\Phi_p(\omega)\). Each infimum is attained at exactly
one repair, so equality holds if and only if that repair is
\(\alpha=0\). For \(p=1\) this says \(b=0\). For \(p=2\) it says
\(a+2b=0\). For \(p=\infty\) it says \(a+b=0\). If the corresponding
linear condition fails, the zero-repair value is strictly larger than
the infimum. \(\blacksquare\)

The lemma is every cochain in \(\mathbb R^2\) and this one line. The
line is not the image of the scalar coupled operator. The lemma does not
compute a quotient distance for a norm other than
\(p\in\{1,2,\infty\}\), a second direction, a third component, or a
vector-valued fiber. The rational content is
\texttt{MinkowskiRepair.skew\_l1\_iff},
\texttt{MinkowskiRepair.skew\_l2\_sq\_iff}, and
\texttt{MinkowskiRepair.skew\_linf\_iff} in
\texttt{lean4/minkowski\_repair.lean}. The symbol \(|2a-b|/\sqrt{5}\)
names the positive square root of the proved square \((2a-b)^2/5\). The
declaration \texttt{quotient\_attained\_minkowski\_chain} still
evaluates a diagonal residual only at \(\alpha=0\). No dashboard margin,
sheaf, or observation map is added.

For the companion discussion, see §7.4.12 of the anchor manuscript, in
the v6.3 release cited as {[}1{]}. False green (Theorem 3) and
irreparable red reflect distinct structural failures: raw inconsistency
in the presence of unobserved directions versus unreachability in the
quotient geometry. Raw section defects belong to \(\operatorname{im}D\);
a positive quotient distance instead identifies a cochain that remains
separated from that fixed repair image. Changing the allowed repair
operator is a different problem.

\subsection{\texorpdfstring{5. What this note does \emph{not}
claim}{5. What this note does not claim}}\label{what-this-note-does-not-claim}

\begin{itemize}
\tightlist
\item
  Nothing here is about any particular AI system, evaluation suite, or
  organization. The theorems constrain the \emph{shape} of composed
  evaluation; whether a real evaluation stack instantiates the setup ---
  and with which kernels --- is an empirical modeling question each
  application must answer and document.
\item
  The results are worst-case and exact, not average-case or statistical.
  They show green scores \emph{can} be uninformative, and exactly when;
  they do not show any particular green dashboard \emph{is}.
\item
  No result here authorizes or forbids any action. A theorem about
  evaluator blind spots is an input to judgment, not a substitute for
  it.
\item
  The severity number \(I(s)\) is one declared diagnostic (a sum of
  interface norms). Other diagnostics make other choices; the companion
  anchor paper {[}1{]} states the axioms under which this additive form
  is forced.
\end{itemize}

\subsection{6. The machine-checked
evidence}\label{the-machine-checked-evidence}

The declaration list below maps the note's results to accompanying Lean
declarations for the abstract observer statements and explicit defect
calculations. The formalization boundaries are described below; in
particular, the graph realization in Theorem 3 is supplied by the
mathematical construction in §3. The recorded evidence, including axiom
inventory (within \texttt{propext}, \texttt{Classical.choice},
\texttt{Quot.sound}) and compile receipts, is in
\texttt{lean4/PROVENANCE.md}. This deposit ships three mathematical Lean
sources under \texttt{lean4/}: \texttt{observer\_visibility.lean}, the
anchor snapshot \texttt{po000\_l004\_observer\_visibility.lean}, and the
mathlib-free rational repair module \texttt{minkowski\_repair.lean}.
\texttt{Census.lean} and \texttt{lakefile.lean} are build
infrastructure; see \texttt{lean4/README.md} and
\texttt{lean4/PROVENANCE.md} for provenance and compile receipts. The
map names this deposit's \texttt{observer\_visibility.lean}. The
byte-identical anchor snapshot
\texttt{po000\_l004\_observer\_visibility.lean} retains the seven core
results under the \texttt{PO000.} prefix; it does not contain either the
joint score-vector strengthening or the Master Soundness Criterion. Both
additions live in this deposit's own source.

All declarations below have the prefix \texttt{ObserverVisibility.}.
Where a subnamespace is specified, append it before the declaration
name.

\textbf{Theorem 1}

\begin{itemize}
\tightlist
\item
  \texttt{tally\_le\_rawNorm}
\item
  \texttt{localGreen\_of\_globalGreen}
\end{itemize}

\textbf{Theorem 2}

\begin{itemize}
\tightlist
\item
  \texttt{allGreen\_iff\_kernel}
\end{itemize}

\textbf{Master Soundness Criterion}

\begin{itemize}
\tightlist
\item
  \texttt{master\_soundness\_iff}
\end{itemize}

\textbf{Theorem 3 --- subnamespace \texttt{ThreePatch}}

\begin{itemize}
\tightlist
\item
  \texttt{three\_patch\_witness}
\item
  \texttt{rawNorm\_scaled}
\item
  \texttt{scaledDefect\_is\_section\_defect}
\end{itemize}

\textbf{Theorem 4 --- subnamespace \texttt{ThreePatch}}

\begin{itemize}
\tightlist
\item
  Maximum-score form: \texttt{no\_recovery\_modulus}
\item
  Full-score-vector form: \texttt{vector\_no\_recovery\_modulus}
\end{itemize}

\textbf{Theorem 5}

\begin{itemize}
\tightlist
\item
  \texttt{rawNorm\_eq\_zero\_of\_allGreen}
\item
  Coverage consequence: \texttt{coverage\_of\_jointTrivial}
\end{itemize}

\textbf{Theorem 6 --- subnamespace \texttt{Coupled}}

\begin{itemize}
\tightlist
\item
  \texttt{full\_coverage}
\item
  \texttt{joint\_kernel\_nontrivial}
\item
  \texttt{allGreen\_forces\_zero}
\end{itemize}

\textbf{Quotient-distance remark --- subnamespace
\texttt{QuotientWitness}}

\begin{itemize}
\tightlist
\item
  \texttt{repair\_mass\_ge\_two}
\item
  \texttt{repair\_mass\_attained}
\item
  \texttt{omega\_not\_exact}
\end{itemize}

\textbf{Repair lemmas --- namespace \texttt{MinkowskiRepair} in
\texttt{minkowski\_repair.lean}}

\begin{itemize}
\tightlist
\item
  Diagonal repair: \texttt{diag\_l1\_iff}, \texttt{diag\_l2\_sq\_iff},
  \texttt{diag\_linf\_iff}.
\item
  The pair \((1,-1)\): \texttt{r2\_l1\_iff}, \texttt{r2\_l2\_sq\_iff},
  \texttt{r2\_linf\_iff}.
\item
  Skew repair: \texttt{skew\_l1\_iff}, \texttt{skew\_l2\_sq\_iff},
  \texttt{skew\_linf\_iff}, \texttt{skew\_pair\_values}.
\item
  Zero-repair, vanishing, and coincident-attainment identities:
  \texttt{diag\_zero\_l1\_iff}, \texttt{diag\_zero\_l2\_iff},
  \texttt{diag\_zero\_linf\_iff}, \texttt{diag\_vanish\_iff},
  \texttt{skew\_vanish\_iff}, \texttt{skew\_points\_coincide\_iff},
  \texttt{skew\_zero\_l1\_iff}, \texttt{skew\_zero\_l2\_iff},
  \texttt{skew\_zero\_linf\_iff}.
\end{itemize}

These 19 reported declarations are over \texttt{Rat}. Euclidean
identities concern squares. The real-valued infimum and square-root
statements in §4 are proved in prose. The finite-dimensional recovery
proposition also has a written proof and no Lean declaration here.

The \texttt{CycleBlindSpot}, \texttt{Redundancy},
\texttt{NonlinearEvaluation}, and \texttt{MinkowskiQuotient}
developments in \texttt{observer\_visibility.lean} are supplementary.
Their historical labels ``Theorem 7'', ``Theorem 7-prime'', and
``Redundancy Corollary'' do not refer to numbered results in this
manuscript, whose numbered theorems are 1--6. They remain available
without enlarging the current paper's claims.

Theorems 1, 2, and 5 are proved over arbitrary finite interface sets,
seminormed state spaces, and contractive measurements. The Master
Soundness Criterion additionally assumes \texttt{hnorm\_zero} on the
supplied realizable set; for arbitrary cochain subsets \(S\) it
establishes \(S \cap K \subseteq \{0\}\), and its
intersection-equals-singleton form \(S \cap K = \{0\}\) uses \(0 \in S\)
(such as \(S = \operatorname{im} D\)). Theorems 3, 4, and 6 are the
explicit finite systems described, with their exact constants. For
Theorem 3, the Lean declarations cover the edge-level defect
calculations described below. The six-site graph realization follows
from the explicit assignment and calculation in §3. The exact Python
replay checks the resulting edge-cochain arithmetic; it does not
independently reconstruct the graph incidence operator.

\emph{Formalization boundary for Theorem 3.} The per-edge realization
declaration listed above formalizes the algebraic identity \(v=v-0\) on
edge-indexed vectors. It does not encode the graph incidence operator or
the six site objects. The graph realization used in this note is
established by the explicit assignment and calculation in §3.

At unit amplitude, the six-interface example has defect mass
\(I(s^{(1)})=6\) and score vector
\((J^0(s^{(1)}), J^1(s^{(1)}), J^2(s^{(1)})) = (0, 0, 0)\). The full
exact-arithmetic replays --- the unit-amplitude assignment, the scaled
family \(I(s^{(\lambda)})=6\lvert\lambda\rvert\), the coupled system of
Theorem 6, and the repair-mass minimum of the Remark --- are in
\texttt{code/lgr\_replay.py}, with every constant pinned by exact
rational equality in \texttt{code/test\_lgr\_replay.py}.

\subsection{7. Provenance}\label{provenance}

This note extracts and restates observer-visibility results from §7.4 of
the author's anchor manuscript, Projection Obstruction for Factorized
States. The corresponding Zenodo release is v6.3 Pre-Print, issued on 4
September 2026, with version DOI
\href{https://doi.org/10.5281/zenodo.22289548}{10.5281/zenodo.22289548}
{[}1{]}. The concept DOI
\href{https://doi.org/10.5281/zenodo.18896776}{10.5281/zenodo.18896776}
identifies the anchor's version family.

This note presents its mathematical statements and proof arguments
directly. The accompanying Lean sources and their recorded provenance
are identified in §6. The distinction there between abstract
formalization and the six-site graph realization applies throughout.

Nothing in this note depends on the anchor's quantum-mechanical or
differential-geometric material.

No external funding. No conflicts of interest.

\subsection{References}\label{references}

{[}1{]} Carroll, J. H. (2026). Projection Obstruction Anchor: Product
Projections and Exact Short-Time Defect Bounds. v6.3 Pre-Print. Zenodo.
https://doi.org/10.5281/zenodo.22289548.

{[}2{]} Boyd, S. (2008). Observability and state estimation. EE263,
Lecture 19, Stanford University, especially slides 19--5 to 19--7 and
19--17. https://ee263.stanford.edu/archive/observ.pdf.

\end{document}
